%============================================================================
% prelude.tex
%   - titlepage
%   - table of contents, list of tables and list of figures
%   - abstract
%============================================================================
\clearpage%
\pagenumbering{roman}  % This makes the page numbers Roman (i, ii, etc)
%--------------------------------------------------------------------%
% TITLE PAGE
%   - define \title{} \author{} \date{}

\title{ Use of  MSPPT in Research Presentation }
\author{Chitra Suresh}

\date{November 2014}

\department{Department of Electronics Engineering}

%    - Set the guide's name
\setguide{Prof. Dr.Vyawahare Vishwesh Anant}
%    - Set the coguide's name (if you have one)
%%\setcoguide{Prof Amitabha Sanyal}

%   - once the above are defined, use \maketitle to generate the titlepage
\maketitle

%--------------------------------------------------------------------
\cleardoublepage \nonumber
%\begin{center}
%\includegraphics{raitlogo.eps}\\
%Ramrao Adik Education Society's\\
%\large {\textbf{Ramrao Adik Institute of Technology}}\\
%\small(Affiliated to the University of Mumbai)\\
%Dr. D. Y. Patil Vidyanagar,Sector 7, Nerul, Navi Mumbai 400
%706.\\\vspace{0.3in} \Large CERTIFICATE
%\end{center}
%\vspace{0.1in}
%\begin{center}
%\large This to certify that Dissertation Seminar-I entitled\\
%\textbf{\lq\lq On Scheduling in WiMAX\rq\rq} \\submitted by
%\begin{center}
% w (09)\\
%\end{center}
%is approved for the degree of \\ \textbf {Master of Engineering}\\ in\\ \textbf {Electronics Engineering}.\\
%
% %\vspace{0.5in}
%%-----------------------------------------------------------------
%%\begin{tabular}{ccc}
%%      \rule{6cm}{1sp}                &\rule{10mm}{0pt}& \rule{6cm}{1sp}
%%      \\
%%       Guide             &&Head of Department \\\\\
%%      \rule{6cm}{1sp}                && \rule{6cm}{1sp} \\
%%      Examiner 1                 && Examiner 2 \\\\\\
%%      Principal: \rule{4cm}{1sp} && \\ && \\
%%       &&
%%    \end{tabular}
%\vspace{1in}
%\begin{tabular}{ccc}
%      \rule{5cm}{1sp}                &\rule{10mm}{0pt}& \rule{5cm}{1sp}
%      \\\vspace{0.5in}
%      Examiner 1                 && Examiner 2 \\
%      \rule{5cm}{1sp}                && \rule{5cm}{1sp} \\ \vspace{0.5in}
%       Guide             && Project Coordinator \\
%      \rule{5cm}{1sp}                && \rule{5cm}{1sp} \\
%      Head of Department                && Principal \\\\
%        %&\rule{5cm}{1sp}& \\
%%       &Principal&
%    \end{tabular}
%\end{center}
%-----------------------------------------------------------------
\cleardoublepage
%--------------------------------------------------------------------%


% ABSTRACT
\begin{abstract}
%\input {abstract}
MS Power Point is a computer program that allows creating and showing slides to support a presentation. The text, graphics and multi-media are combined to content to create research presentations. As a presentation tool MS Power Point can be used to organize, create, structure consistent format of presentation and also provide an illustrative backdrop for the content of presentation and animate slides to give greater visual impact.

Optimization of communication means getting message across and optimization of information under constrains. It can be achieved by means of using less text in each slide.
When presenting ideas that include references to data, it can be helpful to make the point using a graph or table. The data visualization is process which is used to create info graphs.
The spiral chart diagram is power point model used as research process. Setting time for each slide, Hyper links, you tube video are some of the special features of MS Power Point.

Keywords: MS Power Point, optimization of communication, and Info graphs



\end{abstract}

%--------------------------------------------------------------------%
% CONTENTS, TABLES, FIGURES
%\tableofcontents
\listoffigures\

%--------------------------------------------------------------------%
%  Single counter for theorems and theorem-like environments:
\newtheorem{theorem}{Theorem}[chapter]
\newtheorem{assertion}[theorem]{Assertion}
\newtheorem{claim}[theorem]{Claim}
\newtheorem{conjecture}[theorem]{Conjecture}
\newtheorem{corollary}[theorem]{Corollary}
\newtheorem{definition}[theorem]{Definition}
\newtheorem{example}[theorem]{Example}
\newtheorem{figger}[theorem]{Figure}
\newtheorem{lemma}[theorem]{Lemma}
\newtheorem{prop}[theorem]{Proposition}
\newtheorem{remark}[theorem]{Remark}

%--------------------------------------------------------------------%
% Make the page numbers Arabic (1, 2, etc)
\cleardoublepage%
\pagenumbering{arabic}
